\chapter{The Order Book Builder}\label{ch:ll:the-order-book-builder}

On a small-tick instrument, a book kept in balanced trees takes about \qty{70}{\nano\second} per update on this book's
laptop; an array of price levels around the touch takes about \qty{30}{\nano\second}. The array wins, until the day the
stock moves thirty percent: an array of 256 ticks must then be moved and refilled almost twenty thousand times, and the
day's book costs \qty{0.7}{\second} of processor time instead of four milliseconds. The \omterm{def:ll:the-order-book-builder:builder}{book builder} turns the \omterm{def:ll:the-feed-handler:handler}{feed
handler}'s events into the answer every strategy asks first, what is the best price and how much is there, and the answer
must be right, fast and still fast on the day that matters. This chapter states what the book must answer, compares the
structures that answer it by feed type and tick regime, builds an order-by-order book on an open-addressing order map and a
recentring price ladder, and tests it against two books written earlier in the series.

\section{What the book must answer, and how fast}

\begin{definition}[Book builder]\label{def:ll:the-order-book-builder:builder}
A \emph{book builder}\index{book builder} maintains the limit order book (One Quant Book~10, chapter~1) of an instrument from
the \omterm{def:ll:the-feed-handler:handler}{feed handler}'s events: from a market-by-order feed (Book~1, chapter~19) it keeps every resting order and aggregates
them into price levels; from a market-by-price feed it keeps the levels directly. It answers level-1 and level-2 queries
(Book~1, chapter~28) after every event, and knows when its book cannot be trusted.
\end{definition}

A strategy asks three kinds of question. Level~1, the best bid and ask and their sizes, after every event: the microprice of
One Quant Book~7, chapter~8, needs nothing more. Level~2, the first few levels on each side, for depth imbalance and queue
estimates. And, for its own orders, level~3: where is my order in its queue. The builder serves all three from one structure,
at a rate of one update per event of the feed, and the update is on the critical path of chapter~1's budget: a book that
takes \qty{100}{\nano\second} longer per event delays every decision by that much, and at the open, when events come in
bursts (chapter~18), it delays them by the queue of events behind it.

Each event names an order by its reference, so every update starts with a lookup: which order, at what price and side.
Then it changes one price level: adds to it, reduces it, removes it when it empties. An add at a new price creates a level;
a replace (Nasdaq's specification: the original order's ``remaining shares \dots\ must be removed'' and a new reference
``will be used henceforth'') removes an order and adds another. The two costs of a book are therefore the order lookup and
the level update, and the structures below differ in both.

\section{Structures by feed type and tick regime}

\begin{definition}[Dense price ladder]\label{def:ll:the-order-book-builder:ladder}
A \emph{dense price ladder} keeps the price levels of a book in an array indexed by the distance in ticks from a base
price, one slot per tick whether or not a level exists there, covering a window around the touch; levels outside the window
live in a sparse structure, and the window is moved (recentred) when the touch approaches its edge.
\end{definition}

\begin{omfigure}
\begin{tikzpicture}[font=\scriptsize,>=Stealth]
\def\u{0.38}
\foreach \i in {0,...,23} {
  \pgfmathtruncatemacro{\q}{(\i==6)||(\i==8)||(\i==9)||(\i==11)}
  \pgfmathtruncatemacro{\a}{(\i==13)||(\i==14)||(\i==16)||(\i==19)}
  \ifnum\q=1 \draw[fill=omDef!40] (\u*\i,0) rectangle ++(\u,0.5);
  \else\ifnum\a=1 \draw[fill=omThm!35] (\u*\i,0) rectangle ++(\u,0.5);
  \else \draw[fill=gray!6,draw=gray!50] (\u*\i,0) rectangle ++(\u,0.5); \fi\fi
}
\fill[gray!25,opacity=0.5] (0,0) rectangle (\u*6,0.5);
\fill[gray!25,opacity=0.5] (\u*18,0) rectangle (\u*24,0.5);
\draw[decorate,decoration={brace,amplitude=4pt},thick] (\u*6,0.6) -- (\u*18,0.6) node[midway,above=5pt]{middle half: the touch stays here};
\draw[->,omDef,thick] (\u*11.5,-0.7) node[below left=-2pt]{best bid (index 11)} -- (\u*11.5,-0.05);
\draw[->,omThm,thick] (\u*13.5,-0.7) node[below right=-2pt]{best ask (index 13)} -- (\u*13.5,-0.05);
\node[draw,rounded corners=2pt,fill=gray!8,align=center,anchor=west] (far) at (\u*24+0.4,0.25) {sparse map:\\ far levels};
\node[draw,rounded corners=2pt,fill=omMeth!15,align=center,anchor=north west] (map) at (0,-1.6) {order map (open addressing):\\ reference $\to$ slot in the pool};
\node[draw,rounded corners=2pt,fill=omMeth!8,align=center,anchor=west] (pool) at ([xshift=0.8cm]map.east) {order pool:\\ price, quantity, side};
\draw[->] (map) -- (pool);
\node[anchor=north west,align=left] at (0,-2.7) {price $=$ base $+$ index $\times$ tick};
\end{tikzpicture}

\omcaption{The build's book. Levels live in a dense array indexed by ticks from a base price (blue: bids, red: asks, grey:
empty); the best bid and ask are indices kept up to date. When a best price leaves the middle half of the array, the base is
moved to centre the touch and the orders are filed again. Orders are found through an open-addressing map from their
reference to a slot in a fixed pool.}\label{fig:ll:the-order-book-builder:ladder}
\end{omfigure}

Three structures compete for the levels. A balanced tree per side, the textbook answer (\texttt{std::map}, whose search,
removal and insertion are logarithmic and which is ``usually implemented as red-black trees''), allocates a node for each new
level and walks a few of them for each update. A vector of levels sorted best first finds a level by binary search and keeps
the touch at its front, but inserting a level in the middle shifts all the levels behind it. A dense ladder
(\cref{fig:ll:the-order-book-builder:ladder}) finds a level by one subtraction and one division and never allocates; its
costs are elsewhere: the memory of the whole window, the scan to the next non-empty level when the best one empties, and the
recentring when the touch drifts.

The feed type and the tick regime decide. On a large-tick instrument, whose spread is one tick most of the time and whose
levels are all occupied near the touch, every structure is fast, because the book is small and the touch rarely moves: on the
simulator's minute from chapter~18 the three builders take about \qty{28}{\nano\second} (ladder), \qty{36}{\nano\second}
(sorted vector) and \qty{37}{\nano\second} (tree) at the median. On a small-tick instrument, whose orders spread over hundreds
of ticks with gaps between them, the tree and the vector pay for the larger book, 65 to \qty{72}{\nano\second}, while the
ladder stays at about \qty{31}{\nano\second} (\cref{fig:ll:the-order-book-builder:update}). The ladder's tail, about
\qty{150}{\nano\second} at the p99 on the large-tick stream, is its scan: when the best level empties, the next one may be
several empty slots away.

\begin{omfigure}
\begin{tikzpicture}
\begin{axis}[width=12cm,height=5.6cm,ybar,bar width=9pt,ymin=0,ymax=320,
  symbolic x coords={large-tree,large-vector,large-ladder,small-tree,small-vector,small-ladder},xtick=data,
  xticklabels={{tree},{sorted\\ vector},{ladder},{tree},{sorted\\ vector},{ladder}},
  xticklabel style={align=center,font=\footnotesize},enlarge x limits=0.1,ylabel={ns per event},
  tick label style={font=\small},label style={font=\small},grid=major,grid style={gray!25},
  legend columns=2,legend style={font=\footnotesize,at={(0.5,-0.3)},anchor=north,draw=none,fill=none,
  /tikz/every even column/.append style={column sep=8pt}},table/col sep=comma]
\addplot[fill=omDef!70,draw=omDef] table[x=key,y=q50]{figdata/low-latency/19-the-order-book-builder/measured_update_keyed.csv};
\addplot[fill=omThm!45,draw=omThm] table[x=key,y=q99]{figdata/low-latency/19-the-order-book-builder/measured_update_keyed.csv};
\legend{median,p99}
\node[font=\footnotesize,anchor=north] at (axis cs:large-vector,315) {large tick (simulator)};
\node[font=\footnotesize,anchor=north] at (axis cs:small-tree,315) {small tick (synthetic)};
\end{axis}
\end{tikzpicture}

\omcaption{Time to apply one event with three book structures, on the simulator's minute of chapter~18 (a large-tick
instrument, about 190\,000 events) and on 100\,000 synthetic small-tick events (orders up to 300 ticks from the mid). The
timer's own cost is subtracted. Measured on a laptop (Intel Core Ultra 7 155H) under WSL2, no isolated cores. Data:
\texttt{bench\_books.py}.}\label{fig:ll:the-order-book-builder:update}
\end{omfigure}

\section{Order lookup: hash tables and open addressing}

\begin{definition}[Open addressing]\label{def:ll:the-order-book-builder:open}
\emph{Open addressing}\index{open addressing} stores a hash table's entries in one flat array: a key goes to the slot its hash
selects or, if that slot is taken, to the next free one in a fixed probe sequence (linear probing: the following slots), and a
lookup follows the same sequence until it finds the key or an empty slot.
\end{definition}

\texttt{std::unordered\_map}, the standard library's hash table, keeps each entry in a node of its own, allocated on insertion
and reached through a bucket's pointer: one allocation per order and a pointer to follow per lookup. \omterm{def:ll:the-order-book-builder:open}{Open addressing} keeps
the order references and their pool slots in two flat arrays, sized once at start-up to at least twice the largest number of
live orders expected; a lookup is a multiplication, a mask and usually one \omterm{def:ll:memory-hierarchy-and-caches:line}{cache line}. Deletion is the delicate part: emptying
a slot would cut the probe sequences that pass through it, so the build shifts the later entries of the run back into the hole
when their home slot allows it, instead of leaving a tombstone, and lookups stay as short after a day of cancels as at the
open.

\omcode{code/firm/bookbuilder/cpp/firm_bookbuilder.hpp}{37}{53}{Deletion by backward shift: the probe runs stay unbroken
without tombstones.}\label{lst:ll:the-order-book-builder:erase}

The orders themselves live in a pool of fixed-size records allocated at start-up, with a \omterm{def:ll:cpp-for-latency-i-memory:pool}{free list} (chapter~6); the map stores
a slot index, not a pointer. After start-up, applying an event allocates nothing: the build's test counts the heap allocations
while both fixtures are applied, and finds none.

\section{Correctness under gaps, crosses and recovery}

\begin{definition}[Book invariant]\label{def:ll:the-order-book-builder:invariant}
A \emph{book invariant} is a property every correct book satisfies after every event, checked by the builder in its debug
builds: each level's quantity equals the sum of its orders' quantities, no empty level is reported, and the best bid is below
the best ask.
\end{definition}

A book is wrong far more often because of its inputs than because of its structure. An event may name an unknown order (a
reference lost in a gap), reduce an order by more than it holds, or arrive while the handler has flagged the book as stale.
The build's rules follow the \omterm{def:ll:the-feed-handler:handler}{feed handler} of chapter~18: during a snapshot (between its begin and end events) the book is
emptied, rebuilt from the snapshot's orders and flagged stale; an event for an unknown order is ignored and counted; a
reduction larger than the order removes the order. Invariants catch the rest: the Python reference checks them after every
event of both fixtures, and the C++ builder's \texttt{check()} recomputes every level from the pool.

A crossed book deserves a word. The venue never publishes one in continuous trading, so a crossed book in the builder means
a lost event or a bug, and must be treated as stale, not traded on. In auctions and in some consolidated views crosses are
legitimate; the builder then reports them rather than repairing them.

The ladder has one correctness hazard of its own, recentring. When the touch leaves the middle half of the window, the builder
moves the base so that the mid is at the centre and files every live order again: the only pass over all the orders, done in
place.

\omcode{code/firm/bookbuilder/cpp/firm_bookbuilder.hpp}{220}{244}{Recentring: when a best price leaves the middle half, centre
the ladder on the mid and file the live orders again.}\label{lst:ll:the-order-book-builder:recentre}

How often? For a price that moves like a Brownian motion, one expected exit time answers it.

\begin{proposition}[Recentrings of a ladder]\label{prop:ll:the-order-book-builder:recentre}
Let the price, in ticks, be a Brownian motion with standard deviation $\sigma$ per day, and let the ladder be recentred
whenever the price moves $a$ ticks from the centre ($a$ is a quarter of the width). The time between recentrings has mean
$a^2/\sigma^2$ days, so the long-run rate is $\sigma^2/a^2$ recentrings a day. A price that trends by $D$ ticks in a day causes
about $D/a$.
\end{proposition}

\begin{proof}
Let $X_t$ be the move from the centre and $\tau$ the first time $|X_\tau| = a$. Since $X_t^2 - \sigma^2 t$ is a martingale and
$\tau$ has finite mean, optional stopping gives $\E[X_\tau^2] = \sigma^2 \E[\tau]$, and $X_\tau^2 = a^2$, so $\E[\tau] =
a^2/\sigma^2$. The recentrings form a renewal process whose long-run rate is $1/\E[\tau]$. A steady trend travels $a$ ticks
between recentrings, $D/a$ of them in $D$ ticks.
\end{proof}

A stock at 100 with a tick of $10^{-4}$ that moves 1\% a day has $\sigma = 10\,000$ ticks: a ladder of 4\,096 ticks
($a = 1\,024$) recentres about 95 times a day, one of 65\,536 ticks about once every three days. The same stock moving thirty
percent in one day travels 300\,000 ticks: about 290 recentrings with the narrower ladder, and every one of them refiles every
live order. \Cref{fig:ll:the-order-book-builder:width}
measures that day, on synthetic events whose mid drifts thirty percent: 256 recentrings with 4\,096 ticks, fifteen with
65\,536, none with a million, and a total processing time that falls from \qty{0.7}{\second} with a 256-tick ladder to
\qty{4}{\milli\second} at 65\,536 ticks, then rises again as the ladder outgrows the caches.

\begin{omfigure}
\begin{tikzpicture}
\begin{axis}[width=11.5cm,height=5.6cm,xmode=log,ymode=log,log basis x=2,xmin=180,xmax=12000000,ymin=1,ymax=3000,
  xlabel={ladder width (ticks, log scale)},ylabel={ms for the day (log scale)},
  xtick={256,1024,4096,16384,65536,262144,1048576,4194304},
  xticklabels={256,1\,024,4\,096,16\,384,65\,536,$2^{18}$,$2^{20}$,$2^{22}$},
  tick label style={font=\small},label style={font=\small},grid=major,grid style={gray!25},
  nodes near coords,point meta=explicit symbolic,every node near coord/.append style={font=\scriptsize,anchor=south west},
  table/col sep=comma]
\addplot[omDef,very thick,mark=*,mark size=1.6pt] table[x=width,y=total_ms,meta=recentrings]{figdata/low-latency/19-the-order-book-builder/measured_width.csv};
\end{axis}
\end{tikzpicture}

\omcaption{The day a small-tick stock moves thirty percent (100\,000 synthetic events, the mid drifting 300\,000 ticks):
processing time of the ladder book by width, each point labelled with its number of recentrings. Measured on a laptop (Intel
Core Ultra 7 155H) under WSL2, no isolated cores. Data: \texttt{bench\_books.py}.}\label{fig:ll:the-order-book-builder:width}
\end{omfigure}

\section{Testing a book}

A \omterm{def:ll:the-order-book-builder:builder}{book builder} is tested against other books. The build's Python reference, a dictionary of orders and one of levels per side
with nothing clever about it, is compared event by event with the book of One Quant Book~1, chapter~28 (\texttt{firm.feed}),
on its sample of order-by-order messages: the same ten levels on each side after every message. It is compared with the book
behind One Quant Book~7's feature engine (\texttt{firm.lobfeat}) on that engine's fixture: the same best prices and sizes after
every message. The C++ and Rust ladders are then compared with the Python reference on two fixtures, the simulator's events
and a synthetic small-tick stream, through a 64-bit hash of the top five levels on each side, chained over every event; the
comparison runs with the ordinary ladder and with a narrow one that recentres and spills orders into its sparse map, because a
structure's rare paths are where its bugs live. And the tree and sorted-vector builders of the benchmark are checked against
the ladder in the same way before any of them is timed.

\section{Tutorial: three books, two regimes}\label{tut:ll:the-order-book-builder:tut}

\begin{tutorial}
\textbf{Goal.} Build the ladder book, prove it against earlier books, and time it against a tree and a sorted vector on a
large-tick and a small-tick instrument. \textbf{End state:}
\cref{fig:ll:the-order-book-builder:update,fig:ll:the-order-book-builder:width}, and green tests in Python, C++ and Rust.
\begin{tutsteps}
\item \textbf{The fixtures.} \texttt{make\_book\_fixtures.py} writes the \omterm{def:ll:the-feed-handler:handler}{feed handler}'s events for the simulator's three seconds
(chapter~18) and 6\,000 synthetic small-tick events, and the level-2 hashes the reference computes over them.
\item \textbf{Differential tests.} The reference against Book~1's \texttt{firm.feed} and Book~7's \texttt{firm.lobfeat}; the C++
and Rust ladders against the reference (hash after every event), with widths of 4\,096 and 256 ticks; the order map's deletion
on keys that collide; the invariants after every event; zero allocations after construction.
\item \textbf{Measure} with \texttt{python bench\_books.py}: the three structures on the simulator's minute and on 100\,000
small-tick events, then the ladder by width on the day the mid drifts thirty percent.
\end{tutsteps}
\textbf{What to change next.} Replace the ladder's scan for the next non-empty level by a bitmap of occupied levels and a
count-trailing-zeros instruction, and measure the p99 on the large-tick stream again.
\end{tutorial}

\section{Build: the book builder}\label{bld:ll:the-order-book-builder:bookbuilder}

\begin{build}
\bfield{Purpose} The book behind every strategy of the trading system: it consumes the \omterm{def:ll:the-feed-handler:handler}{feed handler}'s events (chapter~18) and
serves level 1 and level 2 to the strategy engine (chapter~20) and the risk gate (chapter~22).
\bfield{Interface} Python reference \texttt{firm\_bookbuilder.Book}: \texttt{apply(event)}, \texttt{best()},
\texttt{depth(side, n)}, \texttt{stale}, \texttt{check()}; \texttt{l2\_hash}; \texttt{recentrings},
\texttt{expected\_recentrings}. C++20 \texttt{firm::book::LadderBook(tick, width, max\_orders)}: \texttt{apply},
\texttt{best}, \texttt{depth}, \texttt{check}, \texttt{stale}, \texttt{recentrings}; \texttt{OrderMap}; \texttt{l2\_hash}. Rust
\texttt{firm\_bookbuilder::LadderBook} and \texttt{OrderMap} with the same.
\bfield{Rules} Order lookup through \omterm{def:ll:the-order-book-builder:open}{open addressing} into a fixed pool; levels in a dense ladder with a sparse fallback; the
ladder recentred when a best price leaves its middle half; the book emptied and flagged stale from a snapshot's begin to its
end; no allocation per event after construction unless an order lands outside the ladder.
\bfield{Acceptance tests} \texttt{code/firm/bookbuilder/}: level 2 identical to Book~1's \texttt{firm.feed} after every message
of its sample and to Book~7's \texttt{firm.lobfeat} after every message of its fixture (reference); the C++ and Rust ladders equal
to the reference after every event of both fixtures, with a normal and a narrow ladder; invariants after every event; zero
allocations; the order map's backward-shift deletion; the recentring model against a simulated Brownian price.
\bfield{Stretch} A bitmap of occupied levels for the scan; level~3 with queue positions for the firm's own orders; a
market-by-price mode; recentring in the background into a second ladder.
\end{build}

\begin{omsources}
\item Nasdaq, \emph{TotalView-ITCH 5.0} specification, section 1.4.5 (Order Replace).
\item cppreference.com, \texttt{std::map}.
\item The earlier books' order books: \texttt{firm.feed} (One Quant Book~1, chapter~28) and \texttt{firm.lobfeat} (One Quant
Book~7, chapter~8).
\end{omsources}

\section{Exercises}

\begin{exercise}[$\star$]\label{exo:ll:the-order-book-builder:1}
A ladder of 4\,096 ticks with a tick of 0.01 is centred on 50.00. What is its base price, which index holds 49.73, and where
must the best bid be for the ladder to stay put?
\end{exercise}

\begin{exercise}[$\star$]\label{exo:ll:the-order-book-builder:2}
How much memory does a ladder of 65\,536 ticks take with 16 bytes per level and two sides? Does it fit the second-level cache
of chapter~3?
\end{exercise}

\begin{exercise}[$\star$]\label{exo:ll:the-order-book-builder:3}
An order of 500 at 20.10 is replaced by one of 300 at 20.11. Which events does the book apply, and what happens to the order's
place in its queue?
\end{exercise}

\begin{exercise}[$\star\star$]\label{exo:ll:the-order-book-builder:4}
Keys 1\,024, 2\,048, 3\,072 hash to the same home slot 5 in a table of 16. They are inserted in that order, then 1\,024 is
erased. Show the table before and after the backward shift. What would a tombstone have cost instead?
\end{exercise}

\begin{exercise}[$\star\star$]\label{exo:ll:the-order-book-builder:5}
A stock trades at 40 with a tick of 0.01 and a daily volatility of 2\%. With a ladder of 1\,024 ticks, how many recentrings
per day does \cref{prop:ll:the-order-book-builder:recentre} predict? And with 4\,096?
\end{exercise}

\begin{exercise}[$\star\star$]\label{exo:ll:the-order-book-builder:6}
The builder receives a reduction of 300 for an order that holds 200. What does it do, and what does it tell the monitoring?
\end{exercise}

\begin{exercise}[$\star\star\star$]\label{exo:ll:the-order-book-builder:7}
\emph{Coding.} Add a bitmap of occupied levels to the C++ ladder (one bit per tick and side) and use count-trailing-zeros to
find the next best level; check the fixtures' hashes and measure the large-tick p99.
\end{exercise}

\begin{exercise}[$\star\star\star$]\label{exo:ll:the-order-book-builder:8}
\emph{Find the flaw.} ``We sized the ladder at 512 ticks because it fits in the first-level cache, and our tests on a quiet
day showed no recentring.''
\end{exercise}

\section{Problem: The Stock That Moved Thirty Percent}

\begin{problem}[{Weekend problem --- how wide a ladder}]\label{pb:ll:the-order-book-builder:1}
A small-tick stock trades at 100 with a tick of $10^{-4}$. On an ordinary day its volatility is 1.5\%; on the day of this
problem it moves thirty percent. Use \cref{prop:ll:the-order-book-builder:recentre} and the measurements of
\cref{fig:ll:the-order-book-builder:width} (\texttt{measured\_width.csv}).

\medskip\noindent\textbf{Part I --- The ordinary day.}
\begin{enumerate}
\item What is $\sigma$ in ticks per day?
\item How many recentrings a day does a ladder of 4\,096 ticks expect? Of 65\,536?
\item What width expects one recentring a day?
\item Why is the ordinary day a poor guide to sizing?
\end{enumerate}

\medskip\noindent\textbf{Part II --- The day that matters.}
\begin{enumerate}[resume]
\item How many ticks does the price travel, and how many recentrings does a steady trend cause with 4\,096 ticks?
\item How many did the measurement count, and why fewer?
\item With 256 ticks the measurement counts far more than the trend alone predicts. Why?
\item What does one recentring cost, from the measurements at 4\,096 and 65\,536 ticks?
\end{enumerate}

\medskip\noindent\textbf{Part III --- The trade-off.}
\begin{enumerate}[resume]
\item Which width minimises the day's measured total, and what does it cost?
\item Why does the total rise again beyond it?
\item How much memory does that width take for both sides?
\item What does the median update time say about the width?
\end{enumerate}

\medskip\noindent\textbf{Part IV --- The verdict.}
\begin{enumerate}[resume]
\item State the \emph{named result}: the ladder width against the recentrings of the day and its total cost, and the width that
minimises it.
\item How would you choose the width per instrument?
\item What would recentring in the background, into a second ladder, change?
\item Why must the tests include a narrow ladder?
\item How does a large-tick instrument change the answer?
\item What should the builder report to monitoring about recentrings?
\item What does the day cost with a balanced tree instead?
\item In one sentence: what does a dense ladder buy, and what must it be protected against?
\end{enumerate}
\end{problem}

\section{Interview questions}

\begin{interviewq}[$\star$ \iqroles{developer}]\label{iq:ll:the-order-book-builder:1}
Design an order book for a market-by-order feed. Which operations must be fast, and how fast?
\end{interviewq}

\begin{interviewq}[$\star\star$ \iqroles{developer}]\label{iq:ll:the-order-book-builder:2}
Tree, sorted vector or array of price levels: when would you choose each?
\end{interviewq}

\begin{interviewq}[$\star\star$ \iqroles{developer}]\label{iq:ll:the-order-book-builder:3}
How would you look up orders by reference without allocating? How do you delete from an open-addressing table?
\end{interviewq}

\begin{interviewq}[$\star\star$ \iqroles{developer}]\label{iq:ll:the-order-book-builder:4}
Your book is crossed. What could have happened, and what do you do?
\end{interviewq}

\begin{interviewq}[$\star\star$ \iqroles{developer}]\label{iq:ll:the-order-book-builder:5}
How do you test a \omterm{def:ll:the-order-book-builder:builder}{book builder} so that you trust it on the day that matters?
\end{interviewq}

\begin{interviewq}[$\star\star\star$ \iqroles{developer, researcher}]\label{iq:ll:the-order-book-builder:6}
A dense ladder is fast until the price moves far. Quantify how far, and design around it.
\end{interviewq}
